\documentclass[10pt,a4paper,openany,twoside]{book}

\usepackage{layout_cours}
\setlist{leftmargin=5mm}

\begin{document}
	\chapnumber{9} % numéro du chapitre
	\titre{Fonctions de référence} % titre du chapitre
	\classe{Seconde} % classe
	\entetegal
    \section*{Corrections}

    
\subsubsection*{120 p.193}
\begin{enumerate}[label=$\alph*)$]
    \item $225<x$ donc $S = \oo{225}{+\infty}$
    \item $0\le x \le \frac{400}{9}$ donc $S = \ff{0}{\frac{400}{9}}$
    \item $0\le x \le 10^{24}$ donc $S = \ff{0}{10^{24}}$
\end{enumerate}

\subsubsection*{122 p.193}
$AB = 10$

$BC = x$

D'après le théorème de Pythagore dans le triangle ABC rectangle en B, 

$AC^2 = 10^2 +x^2$

On résout : 
\begin{align*}
    AC \ge \mathbf{26} &\iff AC^2 \ge 26^2 &&\iff 10^2+x^2 \ge 26^2 &&&\iff x^2\ge 26^2-10^2 \\
    &\iff x\ge \sqrt{26^2-10^2}
\end{align*}

Or $26^2-10^2 = 676 - 100 = 576$ et $24^2 = 576$

Donc $x\ge \sqrt{26^2-10^2} \iff x\ge 24$
\end{document}